\section{The x-free reduced inequality}
\label{sec:envelope}

\emph{The main reduced function.}
Put $A=\sigma^2/s^2$, $r=s^2/L^2\in(0,1/4]$, $g_s=\sqrt{r/(2+r)}$ and $G=P\,g_s$.
The leading terms of Lemma~R are
\begin{multline}
N=2Ar+\frac{GA}{\mu(2+r)}+\frac{3\sigma^2}{8}+\frac{b\sigma^2}{8rA}\\
+\frac{G\sigma^2}{16b}+\frac{G\sigma^2}{16A(2+r)},
\end{multline}
\begin{equation}
D=b+\frac{G}{2}+Ar+\frac{GA}{2\mu(2+r)} .
\end{equation}
Here $g\ge\gamma$ holds if $N-\gamma\sigma D\ge0$.
$N$ is of the form $\alpha A+\beta_1/A+\cdots$, so $A$ can be eliminated exactly by the inequality of arithmetic and geometric means.
The result is a function $\phi(\mu,r,b,G)$; $\phi\ge0$ implies $g\ge\gamma$.
For fixed $(\mu,r,b)$ it is concave in $G$, so the two end points $G=0$ and $G=G_{\max}$ suffice.

\begin{lemma}[Lemma Z3$'$, $x$-free form] \proofst
\label{lem:Z3}
For $x\le0.003$ and $\mu\ge2$, $P\le\Pi^*(\mu)$, where
\begin{equation}
\Pi^*(\mu)=\Qb^2\mu\,\frac{1+at_0}{(1+t_0)^2},
\end{equation}
with
\begin{equation}
t_0=\frac{e}{2}(\mu-0.003)\,e^{(\mu-0.003)/2},\qquad a=c_s e^{-\mu/2},
\end{equation}
and $c_s=e^{0.0015}(1+\sqrt e+2/e)=3.3896$ (the value at $x=0.003$).
For every $\mu>0$, $P\le\mu\Qb^2$.
\end{lemma}
The proof is the fractional knapsack argument of the earlier stages, with $x$ removed at its largest value; it is written out in App.~\ref{app:details:Z3}.
The monotonicity step needs $a\le2$, i.e.\ $\mu\ge1.055$; the cells use $\Pi^*$ only for $\mu\ge2$ and $\mu\Qb^2$ below.
Nothing is lost: the $x$-free envelope reproduces the envelope minimised over $x$.
For example $g\sqrt\ell=1.7116$ at $\sigma=10^{-24}$, $1.3109$ at $10^{-48}$, $1.1508$ at $10^{-96}$, $1.0509$ at $10^{-300}$ and $1.0063$ at $10^{-3000}$ \numst.
The envelope approaches one from above, roughly like $1+40/\ell$.

\emph{Picture.} Of all ripples with the same energy frequency, the hardest one is the one that hides as much energy as the knapsack allows just below the blur.
The envelope is the brightness that even this ripple must leave.
